% Generated by scientific-research-assistant. Compile with:
%   pdflatex report && bibtex report && pdflatex report && pdflatex report
% Use xelatex or lualatex instead if titles contain characters pdflatex rejects.
\documentclass[11pt,a4paper]{article}
\usepackage[utf8]{inputenc}
\usepackage[T1]{fontenc}
\usepackage[margin=2.5cm]{geometry}
\usepackage{booktabs,longtable,array}
\usepackage[hidelinks]{hyperref}
\usepackage{url}
\setlength{\parindent}{0pt}
\setlength{\parskip}{0.6em}

\title{Novelty and reproducibility audit of the UPFD fake-news proposal}
\date{2026-10-18}

\begin{document}
\maketitle


\textbf{Report type:} Research Opportunity Report\\
\textbf{Generated:} 2026-10-18

\textbf{Evidence tags:} [V] verified against a cited source or API response; [E] reasoned estimate; [S] speculative suggestion.

{\small
\begin{longtable}{>{\raggedright\arraybackslash}p{0.307\textwidth}>{\raggedright\arraybackslash}p{0.307\textwidth}>{\raggedright\arraybackslash}p{0.307\textwidth}}
\toprule
\textbf{Constraint} & \textbf{Value} & \textbf{Origin} \\
\midrule
\endhead
field & Graph neural networks for fake news detection (UPFD benchmark) & stated by user \\
experience & undergraduate, comfortable with Python & default (not stated by user) \\
hardware & one consumer GPU, 8 GB VRAM (or a free cloud notebook) & default (not stated by user) \\
gpu vram gb & 8 & default (not stated by user) \\
time budget weeks & 8 & default (not stated by user) \\
language & Python & default (not stated by user) \\
difficulty & moderate & default (not stated by user) \\
mode & any & default (not stated by user) \\
\bottomrule
\end{longtable}
}

\section{Summary}

No-go for the proposal as written: the graph-vs-no-graph comparison on UPFD is already published \cite{karn2025structure,dou2021upfd}. The continuation measured on the official split files that a model-free vote over users seen in training reaches 86.0\% on PolitiFact and 90.3\% on GossipCop, and that repeat users carry identical feature vectors, so every per-graph model can use the shortcut. That shared users predict labels is itself already published \cite{jeong2022hgfnd}; what appears open is treating it as a confound of per-graph benchmark results. Conditional go for that narrower question, scoped to PolitiFact for controlled experiments.

\section{Decisions}

Each shortlisted research question gets go, conditional go or no-go. A go requires full-text reading of the closest work and an explicit check for leakage, weak baselines and evaluation flaws; anything short of that is at most a conditional go.

{\small
\begin{longtable}{>{\raggedright\arraybackslash}p{0.184\textwidth}>{\raggedright\arraybackslash}p{0.184\textwidth}>{\raggedright\arraybackslash}p{0.184\textwidth}>{\raggedright\arraybackslash}p{0.184\textwidth}>{\raggedright\arraybackslash}p{0.184\textwidth}}
\toprule
\textbf{Research question} & \textbf{Decision} & \textbf{Novelty confidence} & \textbf{Mode} & \textbf{Main conditions} \\
\midrule
\endhead
As proposed: multi-seed graph vs no-graph re-evaluation of UPFD & NO-GO & very\_low & replication & not reported \\
Revised (twice): do per-graph UPFD results reflect recurring users rather than content or structure? & CONDITIONAL GO & low & extension & Read \cite{rai2025hypergraph}, \cite{rai2025attention} and \cite{bryla2026networks} through the library; if any evaluates per-graph models on unseen users or reports split overlap, reframe as a replication.; Before running models, state the smallest ac \\
Fallback: independent replication of Karn and Jensen with a held-out test set & NO-GO & low & replication & not reported \\
\bottomrule
\end{longtable}
}

\section{What was verified, and how}

{\small
\begin{longtable}{>{\raggedright\arraybackslash}p{0.307\textwidth}>{\raggedright\arraybackslash}p{0.307\textwidth}>{\raggedright\arraybackslash}p{0.307\textwidth}}
\toprule
\textbf{Dimension} & \textbf{Status} & \textbf{Meaning} \\
\midrule
\endhead
Source existence (verify-refs) & verified: 15 & The cited record exists and its title matches. Says nothing about what the paper shows. \\
Claim support & papers read: abstract: 5, full text: 6, not accessible: 3, title only: 1; claims: estimate: 6, speculative: 1, verified: 23 & Verified claims were checked against their source at the read level shown. \\
Novelty confidence & As proposed: multi-seed graph vs no-graph re-evaluation of UPFD: very\_low (search ceiling low); Revised (twice): do per-graph UPFD results reflect recurring users rather than content or structure?: low (search ceiling low); Fallback: independent replication of Karn and Jensen with a held-out test set: low & How far the searches support a gap; never above medium. \\
Reproducibility & nothing executed & Whether results were re-run, not whether code exists. \\
\bottomrule
\end{longtable}
}

\section{How the literature was searched}

\begin{itemize}
\item \textbf{Sources that answered:} openalex, arxiv, crossref, huggingface
\item \textbf{Sources that failed or were skipped:} semantic\_scholar (rate\_limited / HTTP 429 during challenge searches and paper resolution)
\item \textbf{Search strategies:} challenge, citations of UPFD (OpenAlex citing works), web search, full-text reading of arXiv TeX sources, code reading
\item \textbf{Date range:} 2019 onwards (citing works of UPFD); challenge searches unrestricted
\item \textbf{Records retrieved / after deduplication:} 626 / None
\item \textbf{Search date:} 2026-10-18
\item \textbf{Queries:} \texttt{UPFD fake news graph neural network content-only baseline}; \texttt{does propagation structure help fake news detection}; \texttt{propagation graph ablation fake news detection text-only baseline outperforms}; \texttt{rethinking evaluation fake news detection graph neural networks benchmark}; \texttt{user engagement graph necessary fake news detection GNN}; \texttt{FakeNewsNet propagation-based detection reproducibility}; \texttt{simple baseline fake news detection propagation graphs}; \texttt{user overlap train test leakage fake news detection social context}; \texttt{user-disjoint split fake news detection graph}; \texttt{fake news detection shortcut user identity memorization}; \texttt{spurious correlations social context fake news detection evaluation}; \texttt{unseen users generalization propagation-based fake news detection}; \texttt{FakeNewsNet data leakage evaluation bias}
\item 217 and 223 records examined by the two challenge commands; 186 works citing UPFD screened by title.
\item Karn and Jensen 2025 and Xu 2025 were found by web search, not by the challenge command.
\item Full texts read from arXiv TeX sources: UPFD (2104.12259), Karn and Jensen (2512.06638), DAWN (2411.12775), Xu (2512.09974, partly).
\item Code read: safe-graph/GNN-FakeNews gnn\_model/gnn.py, utils/data\_loader.py, utils/eval\_helper.py, hyper\_parameters.md; PyG torch\_geometric/datasets/upfd.py and examples/upfd.py.
\item Continuation (same day): official split files and small feature files fetched by HTTP range requests from the PyG UPFD archives (about 60 MB of 1.6 GB; no BERT features, no full GossipCop download).
\item Continuation: full text read for HGFND (arXiv 2212.12621) and the Maurya et al. PDF; Springer chapters redirected to a login page and stayed unread.
\item Improved challenge (eval 5) with design elements and the UPFD paper as seed: signal close\_prior\_work\_found, 194 papers via citations, 20 flagged as possibly the same design.
\end{itemize}

\section{What has already been done}

{\small
\begin{longtable}{>{\raggedright\arraybackslash}p{0.115\textwidth}>{\raggedright\arraybackslash}p{0.115\textwidth}>{\raggedright\arraybackslash}p{0.115\textwidth}>{\raggedright\arraybackslash}p{0.115\textwidth}>{\raggedright\arraybackslash}p{0.115\textwidth}>{\raggedright\arraybackslash}p{0.115\textwidth}>{\raggedright\arraybackslash}p{0.115\textwidth}>{\raggedright\arraybackslash}p{0.115\textwidth}}
\toprule
\textbf{Paper} & \textbf{Published} & \textbf{Method} & \textbf{Datasets} & \textbf{Findings} & \textbf{Stated limitations} & \textbf{Read} & \textbf{Link} \\
\midrule
\endhead
User Preference-aware Fake News Detection (Dou et al. 2021) & 2021-07-11 & UPFD: news root node plus user nodes (features from users' historical tweets) in a retweet tree; one GNN layer, readout, concatenation with the news embedding & PolitiFact 314 graphs, GossipCop 5,464 graphs (Table 1) & Full text (TeX source): Table 2 already includes news-only baselines (BERT+MLP 71.04 / 85.76 acc, word2vec+MLP 76.47 / 84.61 on POL / GOS) against UPFD 84.62 / 97.23. Section 3.3 ablation '-EXO' removes all edges: accuracy 85.61 vs 81.63 on PolitiFact and 95.47 vs 93.92 on the second dataset (the text labels both sentences 'Politifact'). Five runs, t-test, 20/10/70 split. & No standard deviations reported. The ablation's full model (GCNFN, word2vec, no news concatenation) reports 85.61, above the 80.54 for GCNFN with word2vec in Table 3; the difference is not explained. & full text & \href{https://doi.org/10.1145/3404835.3462990}{link} \\
The Impact of Data Characteristics on GNN Evaluation for Detecting Fake News (Karn \& Jensen 2025) & 2025-12-07 & Five GNNs (GCN, GAT, GraphSAGE, GCNFN, GNNCL) against an MLP that mean-pools all node features; feature shuffling, edge rewiring and synthetic ego-graphs & UPFD PolitiFact and GossipCop, BERT (768-d) features only & Full text (TeX source): stratified 10-fold CV, mean validation accuracy with 95\% CI. Table 2: GossipCop GCNFN 96.1 +/- 1.6 vs MLP 90.8 +/- 3.0; PolitiFact GCNFN 82.7 +/- 5.5 vs MLP 74.3 +/- 7.4. Edge rewiring leaves accuracy above 90\% on GossipCop; feature shuffling collapses GAT/GraphSAGE to chance while GCN and GNNCL keep about 80\% on GossipCop, which the authors attribute to graph size or density. Over 75\% of nodes are one hop from the root. & Preprint, no code found (author portfolio lists it as a preprint). Text and table disagree: MLP on GossipCop 93.87\% in text vs 90.8\% in Table 2; MLP on PolitiFact 80.93\% vs 74.3\%; GAT on GossipCop 95.26\% vs 91.5\%. Reports validation accuracy only (no held-out test). The MLP pools user nodes too, so it does not separate news content from user features. & full text & \href{https://arxiv.org/abs/2512.06638}{link} \\
Revisiting Fake News Detection: Towards Temporality-aware Evaluation by Leveraging Engagement Earliness (Kim et al. 2025) & 2025-02-26 & Temporality-aware evaluation (70/10/20 temporal split, social context split accordingly) and DAWN edge reweighting by engagement earliness & FakeNewsNet PolitiFact (597 articles) and GossipCop (8,763 articles) built from raw FakeNewsNet, not the UPFD release & Full text (arXiv 2411.12775 TeX source): Table 2, average of 5 runs, no standard deviations: BERT content-only 84.2 / 76.4 acc vs GCN 86.4 / 87.4 and GraphSAGE 85.4 / 87.5 on PolitiFact / GossipCop. Figure caption: users tend to engage with either only fake or only real news. & Code at github.com/LeeJunmo/DAWN has no licence; automated check found no seed control. & full text & \href{https://doi.org/10.1145/3701551.3703524}{link} \\
A New DAWN for Fake News Detection: Exploiting Engagement Earliness for Temporality-aware Evaluation (Kim et al. 2026) & 2026-05-07 & Journal version of DAWN (ACM TIST) & not reported & Abstract only: same temporality-aware evaluation claim as the WSDM paper & not reported & abstract & \href{https://doi.org/10.1145/3815193}{link} \\
Enhancing Fake-News Detection with Node-Level Topological Features (Xu 2025) & 2025-12-10 & Appends degree centrality and local clustering coefficient to UPFD node features & UPFD PolitiFact and GossipCop, BERT features & Abstract and full text: PolitiFact macro-F1 0.7753 to 0.8344; on GossipCop the baseline has slightly better F1 (0.9551 vs 0.9451) & No split, seed or number of runs found in the TeX source & full text & \href{https://arxiv.org/abs/2512.09974}{link} \\
DECOR: Degree-Corrected Social Graph Refinement for Fake News Detection (Wu \& Hooi 2023) & 2023-08-04 & Degree-corrected social graph refinement for fake news detection & Two real-world benchmarks (abstract) & Abstract: degrees of news article nodes show distinctive patterns indicative of veracity & not reported & abstract & \href{https://doi.org/10.1145/3580305.3599298}{link} \\
FANG: Leveraging Social Context for Fake News Detection Using Graph Representation (Nguyen et al. 2020) & 2020-10-19 & FANG: graph representation of the social context (users, sources, news) & not reported & Abstract: learned social-context representations improve fake news detection; cited here because user reuse across news is an intended signal in such social-context models & not reported & abstract & \href{https://doi.org/10.1145/3340531.3412046}{link} \\
Beyond Reported Accuracy: A Verifiability-Focused, Multi-Axis Review of Content and Context Based Fake News Detection (Tuan et al. 2026) & 2026-09-11 & Verifiability-focused review of 45 fake news detection studies & not reported & Abstract: reported performance appears shaped as much by data and evaluation design as by architecture; near-perfect results may reflect separability, leakage or favourable evaluation & not reported & abstract & \href{https://doi.org/10.62411/jcta.17476}{link} \\
A critical look at the evaluation of GNNs under heterophily: Are we really making progress? (Platonov et al. 2023) & 2023-02-22 & Critique of heterophily benchmarks & not reported & Abstract: duplicate nodes cause train-test leakage in standard benchmarks; precedent for a data-overlap audit & not reported & abstract & \href{https://arxiv.org/abs/2302.11640}{link} \\
Multi-dataset Analysis of Misinformation Networks: Divergent Graph Structures and Challenges for Fake News Detection (A \& Kasprzyk 2026) & 2026-01-01 & not reported & not reported & Abstract not retrievable programmatically (publisher login redirect). A web-search summary describes a structural comparison of true and false propagation graphs across platforms; not verified against the text. & not reported & not accessible & \href{https://doi.org/10.1007/978-3-032-14163-7_21}{link} \\
Exploring the Efficacy of Attention Mechanisms in GNN-Based Fake News Detection on UPFD Dataset (Rai \& Rai 2025) & 2025-01-01 & not reported & UPFD (title) & Title only; not read & not reported & not accessible & \href{https://doi.org/10.1007/978-981-97-9855-1_7}{link} \\
Leveraging Hypergraph Neural Network for Fake News Detection on UPFD Dataset (Rai \& Rai 2025) & 2025-01-01 & not reported & UPFD (title) & Title only; not read & not reported & not accessible & \href{https://doi.org/10.1007/978-981-96-4722-4_9}{link} \\
Fake news detection: A survey of graph neural network methods (Phan et al. 2023) & 2023-03-23 & Survey of GNN methods for fake news detection & not reported & Cited as evidence the area is surveyed; abstract not read & not reported & title only & \href{https://doi.org/10.1016/j.asoc.2023.110235}{link} \\
Nothing Stands Alone: Relational Fake News Detection with Hypergraph Neural Networks (Jeong et al. 2022) & 2022-12-24 & HGFND: hypergraph neural network over news items, with hyperedges for shared users, shared time windows and shared entities; transductive node classification & UPFD PolitiFact and GossipCop (news and user information from the UPFD release) & Full text (TeX source): user hyperedges connect news items that share a Twitter user ID (2,953 such users in PolitiFact, 15,646 in GossipCop). With user hyperedges alone the model reaches 87.23 +/- 2.32 (PolitiFact) and 97.29 +/- 0.16 (GossipCop) accuracy; all three hyperedge types give 91.13 / 97.46. Baselines trained with 10 seeds, model chosen on validation. The authors conclude that users play a key role in detecting fake news. & Treats shared users as intended signal in a transductive setting and states there is no label leakage because test labels are hidden; does not evaluate on unseen users & full text & \href{https://arxiv.org/abs/2212.12621}{link} \\
Implementation of Weight Adjusting GNN With Differentiable Pooling for User Preference-aware Fake News Detection (Maurya et al. 2024) & 2024-02-09 & GCN and DiffPool-based GNN on UPFD PolitiFact & UPFD PolitiFact only & Full text (Research Square PDF): single reported run; test accuracy 0.82 for GCN & No seeds, no structure-free baseline, no user-overlap or user-disjoint analysis & full text & \href{https://doi.org/10.21203/rs.3.rs-3942141/v1}{link} \\
\bottomrule
\end{longtable}
}

Method, dataset, finding and limitation cells summarise the cited paper's abstract or text; 'not reported' means the information was not found there.

\subsection{Directions that are already well explored}

\begin{itemize}
\item [V] GNN vs news-only or structure-agnostic baselines on UPFD, including edge removal, edge rewiring and feature shuffling \cite{dou2021upfd,karn2025structure}
\item [V] Content-only vs social-graph models under temporality-aware splits on FakeNewsNet \cite{kim2025dawn}
\item [V] Adding explicit topology features to UPFD node features \cite{xu2025topo}
\item [E] New GNN architectures on UPFD (inferred from titles screened among 186 citing works) \cite{rai2025attention,rai2025hypergraph}
\end{itemize}

\section{Ranked research opportunities}

Each opportunity below survived an explicit search for work that already addresses it. None is claimed to be novel: the novelty confidence states how far the search supports the gap, and what it could have missed.

\subsection{1. Revised (twice): do per-graph UPFD results reflect recurring users rather than content or structure?}
\textbf{Type of gap:} Unverified novelty hypothesis\\
\textbf{Recommended mode:} extension\\
\textbf{Revised:} First revision: the graph-vs-no-graph question is answered \cite{karn2025structure,dou2021upfd}. Second revision (this continuation): full-text reading of \cite{jeong2022hgfnd} showed that the predictive power of users shared between news items is already published, as a deliberate transductive signal. What remains open in the sources read is the evaluation question: how much of the accuracy reported for per-graph models on the official UPFD split, and of conclusions such as 'structure adds little', comes from recurring users that the models can recognise through identical feature vectors.

\textbf{Research question.} On the official UPFD split, how much of the accuracy of per-graph GNNs and of same-feature MLPs is explained by test users seen in training, and do published conclusions about news content versus propagation structure hold when that shortcut is controlled (PolitiFact: masking seen users; GossipCop: descriptive only)?

\textbf{Motivation.} Published comparisons on UPFD attribute performance to news content, user preferences or propagation structure \cite{dou2021upfd,karn2025structure}. Every model that pools user-node features, including Karn and Jensen's MLP, can also exploit recurring users, and DAWN observes that users tend to engage with only one class \cite{kim2025dawn}. A review of fake news studies argues reported accuracies are shaped by evaluation design and possible leakage \cite{tuan2026review}.

\textbf{Supporting papers.} \cite{karn2025structure,kim2025dawn,jeong2022hgfnd,tuan2026review,platonov2023hetero}

\textbf{Closest existing work.}
\begin{itemize}
\item \cite{jeong2022hgfnd} [read: full text; overlap: partial]: Uses shared users as a feature (user hyperedges) and reports their predictive power; does not treat them as a confound of per-graph benchmark results or evaluate on unseen users.
\item \cite{karn2025structure} [read: full text; overlap: partial]: Isolates edges and features with rewiring and shuffling; keeps user nodes and random stratified folds in every condition, so the user shortcut is present throughout.
\item \cite{kim2025dawn} [read: full text; overlap: partial]: Temporal split on raw FakeNewsNet; notes users' label consistency but does not measure user overlap between splits.
\item \cite{dou2021upfd} [read: full text; overlap: different]: Defines the benchmark and split; no overlap analysis.
\item \cite{bryla2026networks} [read: not accessible; overlap: unknown]: Not read; title and a search summary suggest structural analysis, not split leakage.
\item \cite{rai2025hypergraph} [read: not accessible; overlap: unknown]: Not read; title suggests a hypergraph model on UPFD, possibly with user hyperedges.
\end{itemize}

\textbf{Evidence for the gap.}
\begin{itemize}
\item [V] Official split files (PyG/UPFD train\_idx/test\_idx; 62 PolitiFact and 1092 GossipCop training graphs): a vote over the training labels of a test graph's users, with no model, text or edges, reaches 86.0\% on PolitiFact and 90.2\% on GossipCop (majority class 48.9\% and 49.9\%). 79.6\% and 98.8\% of test graphs share a user with training. \cite{dou2021upfd}
\item [V] The shortcut is visible to any model through its inputs: 94.3\% (spaCy history features) and 94.7\% (profile features) of PolitiFact repeat users, and 86.4\% of GossipCop repeat users (profile features), carry an identical feature vector in every graph they appear in. Using the feature vector itself as the identity gives the same vote accuracy (86.0\% PolitiFact, 90.8\% GossipCop). BERT features were not downloaded and not checked. \cite{dou2021upfd}
\item [V] Published accuracies on the same data: UPFD 84.62 (PolitiFact) and 97.23 (GossipCop), its news-only MLPs 71.04-76.47 and 84.61-85.76; the same-feature MLP of Karn and Jensen 74.3 and 90.8; HGFND with user hyperedges only 87.23 and 97.29. The model-free user vote is at or above several of these. \cite{dou2021upfd,karn2025structure,jeong2022hgfnd}
\item [V] That shared users predict veracity is already published: HGFND builds one hyperedge per user shared between news items (2,953 PolitiFact, 15,646 GossipCop, the same counts measured here) and reaches 87.23 / 97.29 accuracy from them alone. It treats this as a legitimate transductive signal, not as a confound of per-graph evaluations. \cite{jeong2022hgfnd}
\item [V] No source read in full (UPFD, Karn and Jensen, DAWN, HGFND, Xu, Maurya) reports train/test user overlap for per-graph models or evaluates on unseen users. \cite{dou2021upfd,karn2025structure,kim2025dawn,jeong2022hgfnd,xu2025topo,maurya2024gnndp}
\end{itemize}

\textbf{Potential contribution.} A short benchmark-validity study: the identity shortcut quantified on the official split, its availability through feature vectors, and whether per-graph models and the 'structure adds little' conclusion change when it is controlled on PolitiFact. Either outcome is reportable; the scope suits a workshop or reproducibility track.

\textbf{Novelty confidence:} low (scale: very\_low, low, medium; a search cannot establish more). The observation that users predict labels is published \cite{jeong2022hgfnd}; only its use as an evaluation confound appears open in the papers read. Three UPFD-related chapters could not be read, Semantic Scholar was rate-limited, and user-disjoint evaluation is common in neighbouring areas under other names (cold-start, inductive users).
Attempts to disprove novelty: 10 search(es); signal \texttt{no\_close\_match\_in\_this\_search}. Queries: \texttt{user overlap train test leakage fake news detection social context}; \texttt{user-disjoint split fake news detection graph}; \texttt{fake news detection shortcut user identity memorization}; \texttt{spurious correlations social context fake news detection evaluation}; \texttt{unseen users generalization propagation-based fake news detection}; \texttt{FakeNewsNet data leakage evaluation bias}; \texttt{propagation graph contribution fake news detection GNN ablation}; \texttt{structure-agnostic baseline graph neural network misinformation benchmark}

\textbf{Feasibility.}
\begin{itemize}
\item [V] A strict user-disjoint split is infeasible on both datasets: shared users link 93.0\% of PolitiFact and 99.6\% of GossipCop graphs into one component; even ignoring users in 3+ graphs the largest component holds 85.4\% and 54.6\%.
\item [V] Masking test-graph users seen in training is a small change on PolitiFact (8.9\% of test user nodes removed; median users per test graph 91 -> 78; 2 graphs emptied) but a large distribution shift on GossipCop (71.6\% removed; median 48.0 -> 6.0; 359 graphs emptied).
\item [E] PolitiFact has 221 official test graphs, so a 95\% interval on accuracy is roughly +/- 5-6 points; only large differences between conditions will be detectable.
\item [V] All checks ran on CPU from about 60 MB of index, label and feature files; models would also be CPU-scale.
\end{itemize}

\textbf{Methodological risks.}

{\small
\begin{longtable}{>{\raggedright\arraybackslash}p{0.23\textwidth}>{\raggedright\arraybackslash}p{0.23\textwidth}>{\raggedright\arraybackslash}p{0.23\textwidth}>{\raggedright\arraybackslash}p{0.23\textwidth}}
\toprule
\textbf{Risk} & \textbf{Status} & \textbf{Evidence} & \textbf{Mitigation} \\
\midrule
\endhead
train test overlap & found & [V] Official split files (PyG/UPFD train\_idx/test\_idx; 62 PolitiFact and 1092 GossipCop training graphs): a vote over the training labels of a test graph's users, with no model, text or edges, reaches 86.0\% on PolitiFact and 90.2\% on GossipCop (majority class 48.9\% and 49.9\%). 79.6\% and 98.8\% of test graphs share a user with training. \cite{dou2021upfd} & Mask seen users in PolitiFact test graphs; report the user vote as a baseline everywhere. \\
weak baseline & found & [V] No published UPFD comparison includes a user-identity vote or a graph-size rule; the size rule reaches 74.9\% on GossipCop's official test split (fails on PolitiFact: 48.4\%). & Include user vote, size rule, root-only MLP and same-feature MLP as baselines. \\
protocol mismatch & found & [V] The official GNN-FakeNews script draws a random split instead of the fixed split files used by PyG, selects no model on validation, and differs from the paper on pooling and weight decay. \cite{dou2021upfd} & Use the fixed split files plus repeated random splits; select on validation. \\
metric definition & found & [V] The official evaluation helper averages macro-F1 over batches. \cite{dou2021upfd} & Compute metrics on the whole test set. \\
insufficient runs & found & [V] 221 PolitiFact test graphs and 62 training graphs; published comparisons use 5 runs. \cite{dou2021upfd} & 10+ seeds and repeated splits with paired comparisons; report intervals. \\
distribution shift & found & [V] Masking test-graph users seen in training is a small change on PolitiFact (8.9\% of test user nodes removed; median users per test graph 91 -> 78; 2 graphs emptied) but a large distribution shift on GossipCop (71.6\% removed; median 48.0 -> 6.0; 359 graphs emptied). & Make causal claims only for PolitiFact; GossipCop analysed descriptively. \\
label leakage & checked ok & [V] graph\_labels match FakeNewsNet's fake/real lists for 313 of 314 PolitiFact and all 5,464 GossipCop graphs (one PolitiFact mismatch unexplained). & not reported \\
temporal leakage & not checked & User features come from historical tweets crawled after the news; their timing was not examined. & not reported \\
\bottomrule
\end{longtable}
}

\textbf{Decision: CONDITIONAL GO.} The shortcut is real, large and measured on the official split, and no source read treats it as a confound of per-graph results. But the signal itself is known \cite{jeong2022hgfnd}, the controlled experiment is only clean on PolitiFact (221 test graphs), and three related chapters are unread. Worth doing as a short, scoped benchmark-validity study, not as a full-semester flagship.

\textbf{Conditions before committing.}
\begin{itemize}
\item Read \cite{rai2025hypergraph}, \cite{rai2025attention} and \cite{bryla2026networks} through the library; if any evaluates per-graph models on unseen users or reports split overlap, reframe as a replication.
\item Before running models, state the smallest accuracy change on masked PolitiFact that would matter, and check it exceeds the roughly +/- 5-6 point interval from 221 test graphs; if not, keep the study descriptive.
\item Frame against \cite{jeong2022hgfnd} (users as a feature) and \cite{karn2025structure} (structure ablation), not as a discovery that users matter.
\end{itemize}

\textbf{Unresolved evidence and risks.}
\begin{itemize}
\item BERT user features not checked for identity (only spaCy and profile features were).
\item Whether the user-feature timing leaks post-publication information is unknown.
\item Reviewers may consider the result implied by HGFND.
\end{itemize}

\subsection{Candidates rejected after the novelty challenge}

\begin{itemize}
\item \textbf{As proposed: multi-seed graph vs no-graph re-evaluation of UPFD} (Already well studied): Reading the full texts showed the experiment has been done three times: the UPFD paper itself reports news-only baselines and an edge-removal ablation \cite{dou2021upfd}; Karn and Jensen benchmark five GNNs against a same-feature MLP with 10-fold CIs, edge rewiring and feature shuffling on both UPFD datasets \cite{karn2025structure}; and DAWN compares a BERT content-only model with GNNs under temporal splits \cite{kim2025dawn}. A further multi-seed rerun would mainly confirm Karn and Jensen. The audit also found the comparison as designed cannot isolate graph structure (see the revised question).
\item \textbf{Fallback: independent replication of Karn and Jensen with a held-out test set} (Limitation reported in prior work): Folded into the revised question: replicating Karn and Jensen without controlling recurring users would repeat the confound; their comparison is re-run inside the PolitiFact masking study instead.
\end{itemize}

\section{Feasibility ranking}

Each criterion is scored 0-5 with evidence. score\_percent is the weighted share earned on scored criteria only; score\_range shows the total if unscored criteria were 0 or 5. Any blocker removes a candidate from the ranking no matter how high its other scores are. Novelty is capped at 2 unless the closest prior work was read in full, and at 1 when that work already ran the experiment.

{\small
\begin{longtable}{>{\raggedright\arraybackslash}p{0.131\textwidth}>{\raggedright\arraybackslash}p{0.131\textwidth}>{\raggedright\arraybackslash}p{0.131\textwidth}>{\raggedright\arraybackslash}p{0.131\textwidth}>{\raggedright\arraybackslash}p{0.131\textwidth}>{\raggedright\arraybackslash}p{0.131\textwidth}>{\raggedright\arraybackslash}p{0.131\textwidth}}
\toprule
\textbf{Rank} & \textbf{Candidate} & \textbf{Status} & \textbf{Score \%} & \textbf{Range} & \textbf{Scored weight \%} & \textbf{Verified weight \%} \\
\midrule
\endhead
1 & Revised: recurring users as a confound of per-graph UPFD results & VIABLE & 70.6 & 70.6-70.6 & 100.0 & 54.0 \\
2 & Fallback: replication of Karn and Jensen with held-out test & VIABLE & 62.6 & 62.6-62.6 & 100.0 & 61.0 \\
3 & As proposed: graph vs no-graph re-evaluation (original O1) & VIABLE & 62.0 & 62.0-62.0 & 100.0 & 71.0 \\
\bottomrule
\end{longtable}
}

\textbf{Scores per criterion (0-5).}

{\small
\begin{longtable}{>{\raggedright\arraybackslash}p{0.184\textwidth}>{\raggedright\arraybackslash}p{0.184\textwidth}>{\raggedright\arraybackslash}p{0.184\textwidth}>{\raggedright\arraybackslash}p{0.184\textwidth}>{\raggedright\arraybackslash}p{0.184\textwidth}}
\toprule
\textbf{Criterion} & \textbf{Weight} & \textbf{B} & \textbf{C} & \textbf{A} \\
\midrule
\endhead
scientific relevance & 12 & 4 [E] & 3 [V] & 3 [V] \\
novelty or replication value & 12 & 2 [E] & 2 [V] & 1 [V] \\
paper and code availability & 10 & 4 [V] & 2 [V] & 4 [V] \\
dataset accessibility & 12 & 3 [V] & 3 [V] & 3 [V] \\
compute fit & 12 & 5 [E] & 5 [E] & 5 [E] \\
implementation simplicity & 10 & 4 [V] & 3 [E] & 4 [V] \\
time fit & 10 & 4 [E] & 5 [E] & 5 [E] \\
reproducibility & 8 & 2 [V] & 1 [V] & 1 [V] \\
evaluation clarity & 7 & 3 [V] & 3 [V] & 2 [V] \\
low failure risk & 7 & 4 [V] & 4 [E] & 2 [E] \\
\bottomrule
\end{longtable}
}

\textbf{Revised: recurring users as a confound of per-graph UPFD results}
\begin{itemize}
\item scientific relevance: UPFD is the standard PyG fake-news benchmark; the diagnostic suggests published comparisons may misattribute performance \cite{karn2025structure,tuan2026review}
\item novelty or replication value: Shared-user signal already published \cite{jeong2022hgfnd}; its role as an evaluation confound not found in six full texts; three chapters unread
\item paper and code availability: Same code and loader as A; overlap script written in this audit
\item dataset accessibility: Mappings and labels downloaded and parsed (under 60 MB)
\item compute fit: Same CPU-scale models as A
\item implementation simplicity: Overlap measurement already scripted; masking is a data transform on PyG graphs
\item time fit: About 6 weeks including the one-week go/no-go checks
\item reproducibility: Diagnostic is scripted and deterministic; upstream training code has the issues listed for A
\item evaluation clarity: Official split files read; but only 221 PolitiFact test graphs
\item low failure risk: Descriptive results already in hand; model runs add to them
\item Missing: novelty\_or\_replication\_value: Three UPFD chapters unread
\item Missing: dataset\_accessibility: Feature files not yet checked; data licence not stated
\end{itemize}

\textbf{Fallback: replication of Karn and Jensen with held-out test}
\begin{itemize}
\item scientific relevance: Most recent direct evidence on the question \cite{karn2025structure}
\item novelty or replication value: Unreviewed preprint with inconsistent numbers; replication value moderate, novelty low
\item paper and code availability: Paper open on arXiv; no code released
\item dataset accessibility: PyG loader
\item compute fit: CPU-scale
\item implementation simplicity: Reimplement five GNNs and perturbations from the description
\item time fit: About 5 weeks
\item reproducibility: No code, seeds or hyperparameters beyond batch size and weight decay
\item evaluation clarity: Published table with CIs to compare against, but it conflicts with the text
\item low failure risk: Confirming or failing to confirm is reportable
\item Missing: dataset\_accessibility: Data licence not stated
\end{itemize}

\textbf{As proposed: graph vs no-graph re-evaluation (original O1)}
\begin{itemize}
\item scientific relevance: Benchmark validity question raised by \cite{karn2025structure,tuan2026review}
\item novelty or replication value: Answered by \cite{dou2021upfd} (-EXO ablation, news-only baselines), \cite{karn2025structure} (MLP vs 5 GNNs with CIs) and \cite{kim2025dawn}
\item paper and code availability: github.com/safe-graph/GNN-FakeNews Apache-2.0; PyG loader and example
\item dataset accessibility: PyG loader, docs HTTP 200
\item compute fit: Graphs of 58-131 nodes on average; one-layer GNNs
\item implementation simplicity: PyG example script exists
\item time fit: About 4 weeks
\item reproducibility: Official script ignores fixed split files, has no validation-based selection, and differs from the paper on pooling and weight decay
\item evaluation clarity: Metrics standard, but published numbers come from differing protocols and a batch-averaged F1
\item low failure risk: The likely outcome, structure adds little, is already published
\item Missing: dataset\_accessibility: Data licence not stated
\end{itemize}

\section{Candidate datasets}

{\small
\begin{longtable}{>{\raggedright\arraybackslash}p{0.102\textwidth}>{\raggedright\arraybackslash}p{0.102\textwidth}>{\raggedright\arraybackslash}p{0.102\textwidth}>{\raggedright\arraybackslash}p{0.102\textwidth}>{\raggedright\arraybackslash}p{0.102\textwidth}>{\raggedright\arraybackslash}p{0.102\textwidth}>{\raggedright\arraybackslash}p{0.102\textwidth}>{\raggedright\arraybackslash}p{0.102\textwidth}>{\raggedright\arraybackslash}p{0.102\textwidth}}
\toprule
\textbf{Dataset} & \textbf{Purpose} & \textbf{Size} & \textbf{Format} & \textbf{Licence} & \textbf{Access} & \textbf{Preprocessing} & \textbf{Suitability} & \textbf{Verified via API} \\
\midrule
\endhead
\href{https://pytorch-geometric.readthedocs.io/en/latest/generated/torch_geometric.datasets.UPFD.html}{UPFD (PolitiFact, GossipCop) via PyTorch Geometric} & All three options & PolitiFact 314 graphs / 41,054 nodes; GossipCop 5,464 graphs / 314,262 nodes (paper Table 1, confirmed by node counts in the id mappings); README states about 1.2 GB for the raw release & PyG InMemoryDataset; loader reads fixed train\_idx/val\_idx/test\_idx files & not reported & reachable (HTTP 200); download via the loader from Google Drive & Choose feature type: profile, spacy, bert, content & Small; the PyG loader uses the fixed split while the original training script does not, so results from the two are not directly comparable & yes \\
\href{https://github.com/safe-graph/GNN-FakeNews/tree/main/data}{GNN-FakeNews id mappings (node id to Twitter user id, node id to timestamp)} & User-overlap measurement (revised question) & pol\_id\_twitter\_mapping.pkl 676 KB; gos\_id\_twitter\_mapping.pkl 6.2 MB & Python pickle of dict[int, str] (opcodes checked: no object construction) & Repository licence Apache-2.0 (covers the repository; data terms not stated) & reachable (HTTP 200) & Graphs rebuilt by walking node ids; root order matched the released news lists exactly & Sufficient to measure overlap without downloading features & yes \\
\href{https://github.com/KaiDMML/FakeNewsNet/tree/master/dataset}{FakeNewsNet fake/real id lists} & Labels for the overlap diagnostic & four CSV files, about 44 MB in total & CSV (id, news\_url, title, tweet\_ids) & not reported & reachable (HTTP 200) & Label by file name (fake/real); counts matched UPFD (157 and 2,732 fake) & Labels only & yes \\
\bottomrule
\end{longtable}
}

\section{Recommendation}

\begin{itemize}
\item [E] Drop the original proposal (no-go) and the stand-alone replication of Karn and Jensen (no-go).
\item [E] Conditional go for the recurring-user question as a 4-6 week benchmark-validity study. Its first deliverable needs no training: the overlap measurements in this report, written up with the code. Model runs on masked PolitiFact follow only if the three conditions are met.
\item [S] If the unread chapters turn out to cover it, or a methods-focused project is preferred, switch to the aspect-based sentiment candidate from the opportunity report.
\end{itemize}

\section{Limitations of this report}

\begin{itemize}
\item Novelty of the revised question rests on searches that missed the closest paper on the original question; treat it as unverified until the library check is done.
\item Graph reconstruction assumes node ids are contiguous per graph with the news node first; this matched the released news list order and the paper's node totals, but was not checked against node\_graph\_id.npy.
\item Accounts appearing in thousands of GossipCop graphs were not identified; removing users in 50+ graphs lowered the GossipCop vote accuracy from 91.1\% to 83.9\%.
\item Karn and Jensen's text and table disagree; the audit quotes the table where it gives CIs and notes the conflict.
\item Semantic Scholar was rate-limited throughout; citing works came from OpenAlex and were screened by title only.
\item No model was trained and no published number was reproduced.
\item Not read: Bryla et al. 2026, Rai et al. 2025 (two chapters); publisher pages require login.
\item BERT user features (125 MB PolitiFact, 862 MB GossipCop) were not downloaded; identity was checked on spaCy and profile features (PolitiFact) and profile features (GossipCop).
\item The user-vote diagnostic shows that identity carries label information; it does not show how much any trained model relies on it. No model was trained.
\end{itemize}

\section{References}

\bibliographystyle{plain}
\bibliography{references}

\end{document}
